\documentclass[../Main.tex]{subfiles}

\begin{document}

\section{Schlussfolgerungen und weiteres Vorgehen} \label{sec:Schlussfolgerungen}

Damit sind wir am Schluss unserer Arbeit angekommen. In diesem Kapitel versuchen wir, unsere Erkenntnisfrage zu beantworten, und gehen dabei auf die Limitationen ein, die wir hatten. Zudem liefern wir weitere Ansätze zur Beantwortung unserer Erkenntnisfrage und schlagen Verbesserungen für unseren Prototyp vor.

\subsection{Beantwortung der Erkenntnisfrage} \label{sec:Beantwortung}

Unsere Erkenntnisfrage lautete:

\begin{quote}
    \textit{Wie müssen Gegnerverhalten und Gegnerplatzierung in einem 3D-Hack-and-Slash-Prototyp konzipiert sein, damit ein zusammenhängender Kampfabschnitt mit nachvollziehbarem Spielfluss, angemessener Schwierigkeit und befriedigendem Spielgefühl entsteht?}
\end{quote}

Für das Gegnerverhalten haben wir bei der Begrenzung der gleichzeitig aktiven Gegner (V2) und bei der Kombination von Nah- und Fernkämpfern (V3) deutliche Hinweise darauf, dass sie so wirken wie beabsichtigt. Dadurch werden nämlich die grossen Gegnermengen, die es in 3D-Hack-and-Slash-Spielen gibt, bekämpfbar, und es entsteht eine Schwierigkeit, die taktischer Natur ist und nicht rein mechanisch durch mehr Lebenspunkte und Schaden. Etwas weniger stark sind unsere Hinweise bei der grundsätzlichen Überlegenheit des Spielers und dem gezielten Durchbrechen dieser Überlegenheit, um Spannung zu erzeugen (V1). Das liegt daran, dass das Durchbrechen nicht immer beabsichtigt war und nicht immer Spannung, sondern auch Frust erzeugt hat. Auch der Vorschlag zur Ankündigung von Angriffen (V4) liefert Hinweise darauf, dass er funktioniert, und zwar um die Nachvollziehbarkeit und das Spielgefühl zu verbessern. Dabei haben wir Hinweise darauf, dass bei den Nahkämpfern Animationen allein ausreichen, bei den Archers aber nicht. Ob eine Ankündigung über weitere Kanäle wie Sound bei den Archers ausreicht, konnten wir demnach nicht überprüfen. Wir würden beide Vorschläge weiterverfolgen, da unsere Daten mindestens einen ausreichenden Hinweis darauf geben.

Für die Gegnerplatzierung hat besonders V8 deutliche Hinweise geliefert, dass neue gegner zuerst alleine im raum platziert werden müssen, bevor sie mit anderen kombiniet werdne. \textcolor{red}{Dass neue Gegnertypen zuerst einzeln und erst danach in Kombination auftreten, haben alle Testpersonen bemerkt, und alle erkannten die neuen Gegner.} Dass die anfangsaufstellung zuerst symetrisch und danach die bewegung der gegner dynamisch ist lieferte auch hinweise, wenn auch weniger starke. wir würden in einem nächsten schritt das dynamische verhalten weg lassen und nur die symetrische platzierung überprüfen. das haben wir dieses mal nicht ganz geschafft weil dieses design zwei unterschiedliche sachen versucht hat zu überprüfen. die grösse des raumes lieferte etwas schwächere hinweise und die aussagen der spieler gehen auseinander und wiedersprechen sich eetwas, deswegen würden wir diesen vorschlag nicht weiter vervolgen. die wellenstruktur  konnten wir tielweise auch beobachten, dass es was bewirkt hat. \textcolor{red}{Bei der lesbaren Anfangsaufstellung mit anschliessend dynamischem Verhalten (V5) wirkte vor allem das Umzingeln, während sich die Lesbarkeit zu Beginn wegen der Doppelfrage Q5 kaum beurteilen liess. Der Wechsel zwischen engen und weiten Räumen (V6) veränderte das Spiel nur bei einem Teil der Testpersonen, und die Wellenstruktur mit Erholungsphasen (V7) bot Erholung vor allem zwischen den Räumen, kaum zwischen den Wellen.} Einige davon würden wir weiterverfolgen. Genauer anschauen würden wir dabei \textcolor{red}{bei V5 die Anfangsaufstellung und das dynamische Verhalten getrennt, bei V6 neben der Grösse auch die Übersichtlichkeit enger Räume und bei V7, wie sich das Ende einer Welle so vermitteln lässt, dass auch zwischen den Wellen Erholung möglich ist.}

durch all diese designvorschläge haben wir deutliche hinweise darauf dass unsere schwierigkeit angemessen war. dass der spielfluss nnachvolziebar war und das spielgefühl gut konnten wir nur teilweise bestätigen, da beide durch faktoren, die nicht dierekt mit den gegnern zu tun hatten beeinflusst wurden, wie z.b. dass das ende eines encounters nicht kommuniziert wurde oder das die combos etwas zu träge waren.

Deshalb lautet die antwort für unsere erkenntnisfrage vorläufig: für unseren prototyp deuten die palytest darauf, dass Gegnerverhalten und Gegnerplatzierung vor allem dann zu einem nachvollziehbaren Spielfluss, einer angemessenen Schwierigkeit und einem befriedigenden Spielgefühl beitragen, wenn die zahl der gleichzietig aktiven gegner bregrenzt, nah und fernkämpfer kobiniert und neue gegner isoliert eingeführt und platziert werden. Die wirkung von wellenstruktur und symetrie hingegen konnten wir nur in ansätzen verbinden.

\textcolor{red}{Bezogen auf die drei Qualitäten der Erkenntnisfrage haben wir für eine angemessene Schwierigkeit deutliche Hinweise, für einen nachvollziehbaren Spielfluss und ein befriedigendes Spielgefühl dagegen nur teilweise: Der Spielfluss wurde vor allem beim Übergang zwischen den Räumen gebrochen, und das Spielgefühl litt unter der Trägheit der Combos, die nicht vom Gegnerverhalten oder der Gegnerplatzierung abhängt (Abschnitt~\ref{sec:Gesamteindruck}).}

\textcolor{red}{Auf die Erkenntnisfrage antworten wir deshalb vorläufig: Für unseren Prototyp deuten die Playtests darauf hin, dass Gegnerverhalten und Gegnerplatzierung vor allem dann zu einem nachvollziehbaren Spielfluss, einer angemessenen Schwierigkeit und einem befriedigenden Spielgefühl beitragen, wenn die Zahl der gleichzeitig aktiven Gegner begrenzt ist, Nah- und Fernkämpfer kombiniert werden und neue Gegner gestaffelt eingeführt werden. Die Wirkung von Raum und Platzierung hing dagegen stärker vom Spielstil der Testpersonen und von der Übersichtlichkeit der Räume ab.}

\begin{itemize}[leftmargin=1.5cm]
    \item \textcolor{red}{Rot = von mir eingesetzt (Bedeutung der Vorschläge, Qualitäten, Antwortsatz), bitte in eigenen Worten umschreiben.}
    \item \textcolor{red}{Noch offen: am Schluss den Bogen zum Einstieg in Kapitel~1 schliessen (Motivation, Spiele genauer verstehen und selbst ein solches Spielerlebnis erschaffen).}
\end{itemize}

\subsection{Limitationen} \label{sec:Limitationen}

Die Ergebnisse dieser Arbeit unterliegen mehreren Einschränkungen.\footnote{Mit Claude formuliert, siehe KI-Deklaration (Abschnitt~\ref{sec:KIDeklaration}).}

Bei den Playtests ist zunächst die kleine Stichprobe zu nennen: Mit fünf Testpersonen lassen sich nur Tendenzen erkennen (Abschnitt~\ref{sec:Auswertung}). Die Urteile in Kapitel~\ref{sec:Evaluation} sind deshalb als Hinweise und nicht als Bestätigungen zu verstehen; wo wir keine Hinweise gefunden haben, bedeutet das nicht, dass das Gegenteil zutrifft. Zudem kannten wir die meisten Testpersonen persönlich oder über Bekannte, was trotz unserer Bitte um ehrliche Antworten zu wohlwollenderen Einschätzungen geführt haben kann (Abschnitt~\ref{sec:AuswahlTP}). Zwei Playtests fanden über Discord statt, wo wir nur Bildschirm und Stimme der Testpersonen wahrnehmen konnten, nicht aber Mimik und Körpersprache (Abschnitt~\ref{sec:Beobachtungsraster}).

Auch das Post-Questionnaire hat Schwächen. Q5 enthält zwei Aussagen auf einmal (überschaubare Anfangsposition und danach unvorhersehbareres Verhalten), sodass unklar bleibt, auf welchen Teil sich eine Bewertung bezieht. Q6 spricht von \emph{dem} engen und \emph{dem} offenen Abschnitt, obwohl der Prototyp mehrere Narrow- und Prospect-Räume enthält (Tabelle~\ref{tab:Raeume}); die Testpersonen könnten dabei unterschiedliche Räume im Kopf gehabt haben.

Bei der Referenzspielanalyse wurde jedes Spiel nur von einer Person analysiert. Die Struktur des MDA-Fragebogens ist zwar einheitlich, Einschätzungen wie \enquote{befriedigend} bleiben aber subjektiv und personenabhängig (Abschnitt~\ref{sec:Analysefragebogen}). Zudem konnten wir die Funktionsweise der Referenzspiele nur aus dem beobachteten Verhalten rekonstruieren, da wir keinen Einblick in deren Quellcode hatten.

\begin{itemize}[leftmargin=1.5cm]
    \item \textcolor{red}{Diesen Abschnitt habe ich aus euren Angaben geschrieben. Falls gewünscht noch ergänzen: mündliche Erklärung der Steuerung vor dem Spielen (statt Selbstentdecken), Gespräche in entspannten Phasen, Leveldesign-Quellen nicht genrespezifisch, Zeitmangel bei der Auswahl der Recherchepunkte.}
\end{itemize}

\subsection{Weiteres Vorgehen und Ausblick} \label{sec:Verbesserungen}

\begin{itemize}[leftmargin=1.5cm]
    \item \textcolor{red}{Optional, als Verbesserungen ausserhalb der Erkenntnisfrage (aus dem gestrichenen Abschnitt \enquote{Weitere Beobachtungen}): Signal am Ende eines Encounters bzw.\ bei offener Tür (TP2, TP4, B10), Combo-Abbruch, Tutorial statt mündlicher Erklärung, Steuerung von Dash und Sprung (TP3, B10; TP2, Q13), Reichweite (TP5, B10), Heilungsdauer (TP4, B10).}
\end{itemize}

\begin{itemize}[leftmargin=1.5cm]
    \item \textcolor{red}{\textbf{Weiteres Vorgehen:} konkrete Verbesserungen am Prototyp, die wir aus den Playtests ableiten (Vertrag C.4: \enquote{Daraus werden konkrete Verbesserungen abgeleitet, die nicht mehr umgesetzt werden müssen}). Das Wort \enquote{Verbesserungen} im ersten Satz verwenden. Jede Verbesserung mit Beleg aus Kapitel~\ref{sec:Evaluation}, z.\,B. akustische Angriffsankündigung, Signal beim Öffnen der Tür, abbrechbare Combos, Tutorial.}
    \item \textcolor{red}{\textbf{Ausblick:} 2--3 Sätze, was man über den Prototyp hinaus untersuchen könnte, z.\,B. mehr Testpersonen, zweiter Playtest nach den Verbesserungen, Spawn-Art als eigene Variable, nicht getestete Prinzipien aus Kapitel~\ref{sec:Leveldesign} (Höhen und Flanken).}
\end{itemize}

\end{document}
